\section{Introduction}

Since its introduction in 1983 as a small name system on the ARPANET, 
the Domain Name System (DNS) has grown to 392.5 million registrations 
across all top-level domains (TLDs)~\cite{dnib_q1_2026}. 
The OpenINTEL project\footnote{\url{https://openintel.nl/}} actively records daily historical data 
of most TLDs. This historical data has been essential for research in many domains 
such as cybersecurity~\cite{Sommese_2024} and infrastructure resilience~\cite{10.1145/3517745.3561458}.
The original design of the OpenINTEL infrastructure~\cite{7460220} stores all data in the 
Parquet\footnote{\url{https://parquet.apache.org/}} columnar storage format for each day and TLD. \\

This research is motivated by the hypothesis that a majority of the collected DNS records 
remain unchanged between measurements. Under this hypothesis a majority of the data stored 
every day by the OpenINTEL project is a duplicate of the previous day. Instead of storing 
all data every day, using a delta-based storage solution, which only stores the records 
that changed compared to a previous snapshot, could substantially reduce the storage usage. 
Furthermore, as the size of the stored records decreases, less data needs to be retrieved 
from storage which could improve the analysis efficiency. The downside of using delta-based 
storage however, is that reconstructing a data point requires reconstructing the full 
dependency chain, which may decrease the analysis efficiency.

Designing such a delta-based storage system presents several non-trivial challenges.\\
Since the data is stored in a format where row-group coherence matters, where it is important 
which row and column a value belongs to, it is not possible to store just the new values of each 
column as that results in not knowing which rows these values might belong to. 
This constrains the system to copy either the full column or the full row into the delta. 
A full-column delta however is impractical for this dataset: a single row receiving a new value would 
require storing all millions of rows in the delta column, yielding no storage benefit. A full-row 
delta follows as the natural choice, but this introduces further problems, namely volatile columns.
The OpenINTEL database contains several volatile columns that change almost every 
measurement, such as the measurement timestamp and round-trip-time. In a row-delta implementation this 
would cause almost all rows to be copied in full every day resulting in no improvement in 
data storage efficiency. However, in the case of a column-delta implementation, a singular row 
receiving a new value would result in all millions of rows being stored in the delta, again resulting 
in almost no storage benefit. From this follows that the system needs a decomposition of the database 
schema to be able to handle both the volatile data while also being able to use delta-compression 
for the non-volatile data.\\

Finally, multiple strategies can be used for a delta-based storage system. For example, 
progressive deltas store deltas based on the previous measurement while snapshot deltas 
would store all deltas based on a single snapshot. These different strategies present tradeoffs
between storage efficiency and reconstruction cost. The exact causes of these tradeoffs lie in 
the characteristics of the underlying data of the OPENIntel system. \\ 

The main focus of this research is the impact of using a delta-based storage system, and its different 
strategy designs on storage and analysis efficiency. This leads to the following research question:

\begin{center}
    \textit{\textbf{
    To what extent can storage and analysis efficiency be improved by transforming a daily historical DNS 
record database into a delta-based storage system?}}
\end{center}

To answer the research question, this paper will first investigate the state of the art in delta-based systems 
to identify the gaps that motivate this work. Furthermore, the paper will evaluate the correctness of the underlying 
hypothesis. Based on the explored system constrains and research gaps, DeltaParq is introduced, a schema-aware delta-based 
storage prototype built for the OPENIntel project. An evaluation framework will be presented that assesses both storage 
and analysis efficiency. With this framework multiple delta strategies will be assessed of which the results will be presented.
Finally, using the results, conclusions are drawn on the practical viability of a delta-based storage system for a daily DNS 
record database.
